\section{Discussion}
\label{sec:discussion}

\subsection{A unifying picture for contradictory reports}

The central contribution of the framework is conceptual: a single computation carries
two robustness numbers, the min-cut $R(G)$ that governs a targeted attacker and the
percolation threshold $q_c$ that governs random damage, and Theorem~\ref{thm:c3} shows
that these diverge as the redundancy width grows. The literature's contradictory
findings then become statements about different threat models applied to the same
circuit. Table~\ref{tab:reconcile} maps the principal empirical observations onto the
framework.

\begin{table}[ht]
\centering
\caption{The framework reads the field's contradictory findings as consistent
statements about distinct robustness numbers of the same circuit.}
\label{tab:reconcile}
\renewcommand{\arraystretch}{1.2}
\begin{tabular}{p{0.46\linewidth}p{0.46\linewidth}}
\toprule
\textbf{Empirical finding} & \textbf{Framework reading} \\
\midrule
Backup name-mover heads: ablating three heads drops the logit difference by
$5\%$ \cite{wang2022ioi} &
Large random robustness; surviving vertex-disjoint paths
(Theorems~\ref{thm:c1}--\ref{thm:c2}). \\
``Redundant paths survived ablation'' \cite{joad2026directions}; refusal cones of
dimension $\le 5$ \cite{wollschlager2025cones} &
$k$-redundancy with $R(G)>1$ (Theorems~\ref{thm:a2},~\ref{thm:c1}); latent and
pre-existing (Proposition~\ref{prop:d1}). \\
Self-repair restoring $\approx 70\%$ of the logit reduction \cite{mcgrath2023hydra} &
The surviving pathways already held the information
(Proposition~\ref{prop:d1}, Theorem~\ref{thm:a1}). \\
Safety removable in $\approx 5$ steps \cite{qi2023finetuning}; shallow and
first-token \cite{qi2024shallow} &
Small effective $R(G)$, hence a tiny targeted budget
(Corollary~\ref{cor:e1}). \\
A single refusal direction \cite{arditi2024refusal} &
The special case $R(G)=1$: the min-cut is the one direction
(Example~\ref{ex:shallow}). \\
\bottomrule
\end{tabular}
\end{table}

\subsection{Relation to the compensatory effect}

The logit-level compensatory effect of \cite{mcgrath2023hydra} is layer-granular and,
as \cite{rushing2024selfrepair} document, roughly $30\%$ of it is a LayerNorm-scaling
artifact rather than genuine rerouting. Theorem~\ref{thm:a1} replaces it with the exact
identity $\Delta I_i=\II(S;R_i\mid R_{-i})$, which by construction discards the
information $R_i$ shares with the other pathways and so carries no scaling artifact.
The result is a head- or feature-level, mechanism-corrected measure of how much a
pathway genuinely contributes that no other pathway can supply.

\subsection{The size of the effect}

The targeted-versus-random separation is not marginal. At depth $d=8$ and width
$w=16$, a random adversary must delete about $110$ of the $128$ internal units---$86\%$
of them---to match the damage a targeted adversary inflicts by removing the $16$
min-cut vertices, $12.5\%$ of the network (Table~\ref{tab:c3}). This is a roughly
sevenfold gap in raw count and a qualitative gap in fraction ($q_c=0.86$ against the
targeted $1/d=0.125$). Redundancy buys a great deal of robustness against
undirected damage and almost none against an adversary who attacks the cut.

\subsection{Limitations}

We state the framework's boundaries honestly.

\begin{itemize}[leftmargin=*,itemsep=2pt,topsep=2pt]
  \item \emph{Combinatorial idealization.} The graph model treats $S$ as computable
    only if a path is fully intact. Real transformers are analog: sub-threshold
    contributions along several partially damaged paths can sum past a decision
    threshold, so the min-cut $R(G)$ can \emph{under}-state true robustness. The exact
    survival law of Theorem~\ref{thm:c2} also assumes complete bipartite wiring between
    bundles; sparser wiring changes the constants but not the qualitative transition.
  \item \emph{Estimation on real models.} Computing $\Delta I_i=\II(S;R_i\mid R_{-i})$
    requires estimating conditional mutual information in high dimensions, which is
    hard and biased; the partial-information cross-check used here can over-report
    redundancy. We did not instantiate the bounds on an actual LLM---this is a theory
    plus synthetic-verification contribution, not an empirical measurement.
  \item \emph{Fano gives a floor.} Theorem~\ref{thm:b1} lower-bounds the behavioral
    error; it does not assert that the floor is achieved (it is tight when $S\mid
    R_{T^c}$ is uniform on its support).
  \item \emph{Static $R(G)$.} Corollary~\ref{cor:e1} bounds the budget to misalign at
    a fixed network but does not model how $R(G)$ itself evolves under training.
  \item \emph{Mapping behavior to a variable.} Real aligned behavior is not a single
    discrete variable; the framework applies per chosen behavior $S$.
\end{itemize}

\subsection{Open questions}

\begin{itemize}[leftmargin=*,itemsep=2pt,topsep=2pt]
  \item \emph{Dynamics of $R(G)$ under training.} Determine conditions under which
    adversarial or continued training grows the alignment min-cut rather than cutting
    it. The capacity bound of \cite{hanni2024superposition} and the
    abundance/scarcity regimes of \cite{bereska2025superposition} are candidate
    predictors.
  \item \emph{Analog-threshold percolation.} Replace ``path intact'' by ``summed
    contribution exceeds a threshold $\theta$'' and derive the corresponding weighted
    min-cut and flow threshold, lifting the combinatorial idealization.
  \item \emph{Degeneracy beyond the single min-cut.} Port the networked-buffering and
    degeneracy notions of \cite{whitacre2010degeneracy} to bound multi-pathway lesion
    robustness by a graph quantity richer than the single min-cut.
  \item \emph{Tight behavioral bounds.} Develop list-decoding analogues of
    Theorem~\ref{thm:b1} for graded refusal behaviors and relate the decoder error
    $P_e$ quantitatively to the output robustness $\rho$ for vector behaviors.
  \item \emph{Estimator-robust ablation.} Design a redundancy-aware estimator of
    $\Delta I_i$ that resists the underestimation of single-head ablation reported at
    scale by \cite{frank2026routing}.
\end{itemize}
